% 13_application_examples.tex
\chapter{Application Examples}
\label{ch:vol9_application_examples}

\begin{objectivebox}
\begin{itemize}
    \item Document a curated set of application examples in which substrate primitives from Chs.~2--12 combine to produce engineering observables: the $\alpha$ closed-form identification at the electron LC tank (Class-B, retained input); Schwarzschild reduction as a SYM-class gradient-index transmission line; $W$/$Z$ gauge-boson masses as transformer-leakage cutoff modes; the Born-rule $p=2$ exponent as substrate boundary-Joule extraction at a $\Gamma$ port; $\delta_{strain}$ as the substrate's cosmic-scale TCC entry.
    \item Frame each example as cross-chapter integration: each application invokes 3--5 substrate primitives across distinct datasheet chapters and resolves them at canonical leaves via \texttt{\textbackslash ref\{\}}. No new substrate-physics is derived here; the chapter is a navigation map.
    \item Preserve the load-bearing framings carried through every preceding chapter: substrate is natural (no engineered applications); $\nu_{vac} = 2/7$ is the Ax~1 LC + Ax~2 $\alpha$ interlink (not promoted); $c_{EM}$ vs $c_{shear}$ are not interchangeable (Pitfall~\#5); Machian $G$ is mixed --- form-derived (the Achromatic-Lens) but value calibrated ($\xi$ calibrated to CODATA $G$; \kbleaf{common/interlock-register.md}), not a primitive axiom; lattice-genesis is single-seed \emph{within our cosmic horizon} (one seed, not many simultaneous disjoint seeds sharing our horizon --- not a claim that ours is the only universe); BH interiors are Regime~IV ruptured plasma (not a B-sublattice), and a melted interior \emph{can} re-crystallize into a daughter cosmology --- nested parent/daughter cosmologies across $\Gamma=-1$ melt$\to$recrystallize boundaries are permitted (Ch.~\ref{ch:vol9_cosmological_characteristics}~$\S$\ref{sec:vol9_cosmo_substrate_rupture}; our own universe is a daughter per \kbleaf{common/omega-freeze-cosmic-grain-cascade.md}).
    \item Tag each example's discipline class per \texttt{ave-discrimination-check} + \texttt{consistency-vs-emergence} v1.3: substrate-distinct-prediction-with-named-SM-counterfactual, identity-class consistency, manifestation-class, or substrate-mechanism mapping to observed quantities.
    \item Provide a Cross-Chapter Citation Map for reader navigation: rows are application examples, columns are Vol~9 chapters whose substrate primitives the example invokes.
\end{itemize}
\end{objectivebox}

\section{Cross-Chapter Integration}
\label{sec:vol9_app_cross_chapter_integration}

The preceding chapters of this datasheet catalog the substrate's individual primitive parameters --- the LC bond ($\varepsilon_0$, $\mu_0$, $Z_0$, $c_0$, $\ell_{node}$, $\alpha$, $\xi_{topo}$; Ch.~\ref{ch:vol9_dc_electrical_characteristics}); the AC resonance and operating-point modulation ($\omega_C$, $c_{EM}(A_0)$, $c_{shear}(A_0)$; Ch.~\ref{ch:vol9_ac_electrical_characteristics}); the Cosserat-Curie thermal sector ($\delta_{strain}$ at $T_{CMB}$; Ch.~\ref{ch:vol9_temperature_characteristics}); the Axiom~4 saturation kernel and its four-regime partition (Ch.~\ref{ch:vol9_saturation_characteristics}); the two-threshold breakdown structure ($V_{yield}$, $V_{snap}$, $E_S$; Ch.~\ref{ch:vol9_breakdown_characteristics}); the Cosserat mechanical primitives ($G_{vac}$, $K_{vac} = 2 G_{vac}$, $\nu_{vac} = 2/7$, $\gamma_c$; Ch.~\ref{ch:vol9_mechanical_characteristics}); the microrotational / magnetic sector ($\mu_0$, $L_{cell}$, $l_c$, $m_\omega$; Ch.~\ref{ch:vol9_magnetic_microrotational_characteristics}); the topological invariants ($(2,3)$ Clifford-torus electron, $I4_1 32$ chirality, $0_1$ unknot real-space body; Ch.~\ref{ch:vol9_topological_characteristics}); and the cosmic-scale primitives ($R_H$, $\hat{\Omega}_{freeze}$, $u_0^*$, Machian $G$; Ch.~\ref{ch:vol9_cosmological_characteristics}).

This chapter shows how those primitives combine. Each application example below walks a single observed physical phenomenon back to the substrate primitives that produce it; each primitive is cited via cross-chapter \texttt{\textbackslash ref\{\}} into the chapter where that primitive's spec table lives, and into the canonical KB leaf where the substrate-physics derivation lives. No derivation is repeated; the chapter is a guided tour of cross-chapter combination.

\textbf{Substrate-natural framing (load-bearing).} Each example below is a substrate-mechanism mapping --- the substrate is what it is; the observed phenomenon is what the substrate naturally produces; AVE substrate-physics names the mechanism; engineering practice empirically measures the consequence. Examples like the $\alpha = 1/(4\pi^3 + \pi^2 + \pi + \text{correction})$ derivation and the Born-rule $p = 2$ exponent are not ``framework-distinct claims'' that engineer a new physics --- they are derivations of standard observed quantities from the substrate's natural primitives, with the SM-language and substrate-language descriptions identified as two coordinate systems for the same physics. \texttt{ave-discrimination-check} discipline applies: where an example carries a substrate-distinct prediction beyond what the SM coordinate alone reaches (e.g., $\delta_{strain}$ as a thermal-mode-population ASYM modulation), the SM counterfactual is named explicitly.

\section{Example: Electron LC Tank $\to$ Fine-Structure Constant}
\label{sec:vol9_app_electron_lc_alpha}

\textbf{Observed phenomenon.} The fine-structure constant $\alpha^{-1} = 137.035999$ (CODATA, 12 decimal places).

\textbf{Substrate primitives invoked.}

\begin{itemize}
    \item \textbf{$(2,3)$ Clifford-torus phase-space winding} for the electron (Ch.~\ref{ch:vol9_topological_characteristics}~$\S$\ref{sec:vol9_torus_knot_winding}); $0_1$ unknot real-space body at minimum ropelength $2\pi$.
    \item \textbf{Bond LC primitives} $L_{cell} = \mu_0\,\ell_{node}$, $C_{cell} = \varepsilon_0\,\ell_{node}$ (Ch.~\ref{ch:vol9_dc_electrical_characteristics}~$\S$\ref{subsec:vol9_dc_impedance_trio}); bond TLM characteristic impedance $Z_0 = \sqrt{\mu_0/\varepsilon_0} \approx 376.73\,\Omega$.
    \item \textbf{Topological transduction} $\xi_{topo} = e/\ell_{node}$ (Ch.~\ref{ch:vol9_dc_electrical_characteristics}~$\S$\ref{subsec:vol9_dc_coupling_primitives}; Ch.~\ref{ch:vol9_topological_characteristics}~$\S$\ref{sec:vol9_topological_spec_table}).
    \item \textbf{Axiom 3 minimum-reflection principle} (the substrate minimizes $|\Gamma|^2$ at every internal boundary; Ch.~\ref{ch:vol9_pin_port_configuration}~$\S$\ref{sec:vol9_substrate_boundary_semantics}); $\Gamma = -1$ TIR closure at the saturated bond (Ch.~\ref{ch:vol9_pin_port_configuration}~$\S$\ref{sec:vol9_boundary_saturation_conditions}).
    \item \textbf{Magic-angle operating point} $u_0^* \approx 0.187$ (value calibrated to CODATA $\alpha$, $G$, not forward-derived; the lock $S(A^*)=0$ / $K=2G$ is firm) at which the Axiom~4 kernel closes ($S(A^*) = 0$) and the Golden-Torus geometry $(R, r, d) = (\varphi/2, (\varphi-1)/2, 1)$ is realized (Ch.~\ref{ch:vol9_cosmological_characteristics}~$\S$\ref{sec:vol9_cosmo_magic_angle}).
\end{itemize}

\textbf{Cross-chapter chain.} The electron is a $(2,3)$ trefoil in Clifford-torus phase-space coordinates on the K4 substrate (Ch.~\ref{ch:vol9_topological_characteristics}~$\S$\ref{sec:vol9_torus_knot_winding}) with real-space body the $0_1$ unknot at minimum ropelength. At the magic-angle operating point $u_0^*$ (Ch.~\ref{ch:vol9_cosmological_characteristics}~$\S$\ref{sec:vol9_cosmo_magic_angle}), the substrate's bond pair saturates at $\Gamma_{\mathrm{bulk}} = -1$ at both ends (Ch.~\ref{ch:vol9_pin_port_configuration}~$\S$\ref{sec:vol9_boundary_saturation_conditions}; device schematic Fig.~\ref{fig:vol9_circuit_electron_barrier}), forming a TIR-confined LC tank whose stored-energy quality factor is set by the codimensional mode-count assembly at the Golden-Torus geometry. The Golden-Torus closure (\kbleaf{ave-kb/vol1/ch8-alpha-golden-torus.md}) is the addition of three independent codimensional Nyquist-cell mode counts at the $\Gamma = -1$ boundary:

\[
\alpha^{-1}_{\text{ideal}} \;=\; \underbrace{4\pi^3}_{\Lambda_{\text{vol}}\;(3\text{-cycle})} + \underbrace{\pi^2}_{\Lambda_{\text{surf}}\;(2\text{-cycle})} + \underbrace{\pi}_{\Lambda_{\text{line}}\;(1\text{-cycle})} \;\approx\; 137.0363038.
\]

The three terms are: the $4\pi$ temporal-phase closure per observable Compton cycle from bipartite K4 lobe-count (2 sublattices) $\times\, 2\pi$ phasor rotation per lobe (substrate-native, \kbleaf{ave-kb/vol2/particle-physics/ch01-topological-matter/l3-electron-soliton-synthesis.md}:103--105; the standard-physics translation reference is the $T \to 2T \subset SU(2)$ K4-rotation-group double-cover applied temporally per Ch.~\ref{ch:vol9_topological_characteristics}~$\S$\ref{sec:vol9_su2_so3_double_cover}); the Clifford-torus 2-area at substrate-derived $R \cdot r = 1/4$ (the phasor enclosed area at Axiom-4 self-saturation onset equals the Nyquist cell cross-section area); and the Ampère 1-cycle around the transverse cross-section perimeter at Nyquist-quantized diameter. The cold-lattice value $\alpha^{-1}_{\text{ideal}} = 4\pi^3 + \pi^2 + \pi$ is the substrate's intrinsic value at $T \to 0$, $A_0 = 0$, $S(0) = 1$ (canonical \texttt{ALPHA\_COLD\_INV}, \kbleaf{src/ave/core/constants.py}). The observed CODATA $\alpha^{-1} = 137.035999$ sits below the cold-lattice asymptote by the fractional vacuum strain $\delta_{strain} \approx \num{2.22e-6}$ documented in \S\ref{sec:vol9_app_delta_strain_cosmic_tcc} below (the Cosserat-Curie thermal-mode-population ASYM correction at $T = T_{CMB}$).

\textbf{Discipline class.} Per \texttt{consistency-vs-emergence} v1.3: the cold-lattice closure $\alpha^{-1}_{\text{ideal}} = 4\pi^3 + \pi^2 + \pi$ is a Class~B substrate-mechanism mapping --- the SM-language fine-structure constant identifies as the substrate's $\Gamma = -1$ TIR-mode count at the electron LC tank's Golden-Torus geometry. The decimal accuracy $|\alpha^{-1}_{\text{CODATA}}/\alpha^{-1}_{\text{ideal}} - 1| \approx \num{2.2e-6}$ is the $\delta_{strain}$ correction, a \textbf{definitional residual} ($1-$CODATA$/\alpha_{cold}$): the substrate-mechanism class predicts the SIGN; the magnitude is not derived from substrate statistical mechanics (a candidate thermal-occupation derivation undershoots by $\sim$31 orders of magnitude and is not substrate-distinct), so the magnitude is a definitional residual (Ch.~\ref{ch:vol9_temperature_characteristics}~$\S$\ref{sec:vol9_temperature_open}).

\section{Example: Gravity as Gradient-Index Transmission Line}
\label{sec:vol9_app_gravity_gradient_index}

\textbf{Observed phenomenon.} Schwarzschild gravitational time dilation $c_{group}(r) = c_0 \sqrt{1 - r_s/r}$ at the weak-field limit, with simultaneously null $\Delta\alpha/\alpha$ between observers at different gravitational potentials.

\textbf{Substrate primitives invoked.}

\begin{itemize}
    \item \textbf{Cosserat bulk moduli} $G_{vac}$, $K_{vac} = 2 G_{vac}$, $\nu_{vac} = 2/7$ (Ch.~\ref{ch:vol9_mechanical_characteristics}~$\S$\ref{sec:vol9_mech_poisson}); the $K/G = 2$ EMT closure at packing fraction $p^* = 8\pi\alpha$ is the substrate-axiomatic interlink (Ax~1 LC + Ax~2 $\alpha$).
    \item \textbf{Two effective wave speeds} $c_{EM}(A_0) = c_0/S(A_0)$ (Maxwell phase velocity) and $c_{shear}(A_0) = c_0\sqrt{S(A_0)}$ (mechanical / group / rest-mass velocity) (Ch.~\ref{ch:vol9_saturation_characteristics}~$\S$\ref{sec:vol9_two_wave_speeds}); Ch.~\ref{ch:vol9_mechanical_characteristics}~$\S$\ref{sec:vol9_mech_wave_speeds}; CLAUDE.md INVARIANT-S2 Pitfall~\#5).
    \item \textbf{SYM-class scaling} (gravity-class realization): $\mu(r)$ and $\varepsilon(r)$ scale together by the same factor so $Z_0 = \sqrt{\mu/\varepsilon}$ stays invariant (\kbleaf{ave-kb/vol3/gravity/ch01-gravity-yield/alpha-invariance-symmetric-gravity.md}) (Ch.~\ref{ch:vol9_pin_port_configuration}~$\S$\ref{sec:vol9_port_class_taxonomy}; Ch.~\ref{ch:vol9_saturation_characteristics}~$\S$\ref{sec:vol9_alpha_invariance}).
    \item \textbf{Machian $G$ as integrated boundary impedance --- mixed (derived form, value-fitted $\xi$)} $G = c^4/(7\xi T_{EM}(u_0^*))$ with dimensionless Machian-hierarchy coupling $\xi = 4\pi(R_H/\ell_{node})\alpha^{-2}$ (form derived; $\xi$ calibrated to CODATA $G$; $u_0^*$ value calibrated; \kbleaf{common/interlock-register.md}, \kbleaf{full-derivation-chain.md}) (Ch.~\ref{ch:vol9_cosmological_characteristics}~$\S$\ref{sec:vol9_cosmo_machian_g}).
    \item \textbf{Axiom 3 minimum-reflection} at every substrate port (Ch.~\ref{ch:vol9_pin_port_configuration}~$\S$\ref{sec:vol9_substrate_boundary_semantics}): SYM-class ports satisfy $\Gamma = 0$ identically at every interior point.
\end{itemize}

\textbf{Cross-chapter chain.} Gravitational field gradients induce SYM-class scaling of the substrate's bulk constitutive parameters: $\mu(r)$ and $\varepsilon(r)$ scale by the same factor across the field region, so $Z_0$ is preserved interior-wide ($\Gamma = 0$ everywhere by Op17 + Axiom~3; Ch.~\ref{ch:vol9_pin_port_configuration}~$\S$\ref{sec:vol9_port_class_taxonomy}). The substrate's two effective wave speeds respond differently to this scaling: $c_{shear}(A_0) = c_0 \sqrt{S(A_0)}$ tracks the Schwarzschild $\sqrt{1 - r_s/r}$ reduction (Ch.~\ref{ch:vol9_mechanical_characteristics}~$\S$\ref{sec:vol9_mech_wave_speeds}), giving observed gravitational time dilation; $c_{EM}(A_0) = c_0/S(A_0)$ enters $\alpha = e^2/(4\pi\varepsilon_0\hbar c)$ in the inverse direction. The asymmetry is exact: under SYM scaling by factor $nS$, $\alpha_{eff}/\alpha_0 = e^2/[4\pi(\varepsilon_0 nS)\hbar(c_0/(nS))] = \alpha_0$. The substrate's gravity-class realization simultaneously produces gravitational time-dilation via $c_{shear}$ AND preserves $\alpha$ exactly (\kbleaf{ave-kb/vol3/gravity/ch01-gravity-yield/alpha-invariance-symmetric-gravity.md}; Ch.~\ref{ch:vol9_saturation_characteristics}~$\S$\ref{sec:vol9_alpha_invariance}).

The Machian impedance integral closes the macroscopic gravitational coupling: a local test mass couples to every other mass within $R_H$ through the substrate's integrated boundary impedance, giving Newton's $G = c^4/(7 \xi T_{EM}(u_0^*))$ with the factor of 7 inherited from $\nu_{vac} = 2/7$ (Ch.~\ref{ch:vol9_mechanical_characteristics}~$\S$\ref{sec:vol9_mech_poisson}; Ch.~\ref{ch:vol9_cosmological_characteristics}~$\S$\ref{sec:vol9_cosmo_machian_g}). The substrate's gravitational coupling \emph{form} is therefore tied to the cosmological operating point $u_0^*$ --- but $G$ is mixed (form-derived, value calibrated to CODATA $G$; \kbleaf{common/interlock-register.md}) and $u_0^*$ is itself calibrated to CODATA $\alpha$, $G$, not forward-derived (Ch.~\ref{ch:vol9_cosmological_characteristics}~$\S$\ref{sec:vol9_cosmo_magic_angle}). The electromagnetic ($\alpha$) and gravitational ($G$) routes fix $u_0^*$; the independent cosmological $\mathcal{J}_{cosmic}$ route is the test of it (pass $=$ chord, fail $=$ echo; Ch.~\ref{ch:vol9_falsification_tests}~$\S$\ref{sec:vol9_three_route_falsifiability}).

\textbf{Discipline class.} The Schwarzschild-reduction substrate mapping is Class~B (substrate-mechanism mapping to a derived observable); the SYM-class $\alpha$-invariance is a Class~B substrate-mechanism mapping (\kbleaf{ave-kb/vol3/gravity/ch01-gravity-yield/alpha-invariance-symmetric-gravity.md}). Per \texttt{ave-discrimination-check}: the SM-counterfactual --- GR gravitational time dilation derived from $G_{\mu\nu} = 8\pi T_{\mu\nu}$ via the Einstein equations --- gives identical $c_{group}(r)$ at the weak-field limit; the substrate version reaches the same result via a different mechanism (gradient-index TL on a Cosserat micropolar continuum), and adds a falsifiable distinction at strong gravitational field where the substrate's underlying Axiom~4 kernel governs all regime-transitions including BH-formation (Ch.~\ref{ch:vol9_cosmological_characteristics}~$\S$\ref{sec:vol9_cosmo_substrate_rupture}: BH interior = Regime~IV ruptured plasma). The Machian-$G$ \emph{form}-derivation is substrate-distinct beyond GR's posited $G$ (though $G$'s value is calibration-fitted, \kbleaf{common/interlock-register.md}) --- the three-route $u_0^*$ falsifier ($\alpha$ and $G$ fix the operating point, the independent $\mathcal{J}_{cosmic}$ route tests it) is documented at Ch.~\ref{ch:vol9_cosmological_characteristics}~$\S$\ref{sec:vol9_cosmo_omega_freeze}.

\section{Example: $W$/$Z$ Gauge Bosons as Transformer-Leakage Cutoff Modes}
\label{sec:vol9_app_wz_transformer_leakage}

\textbf{Observed phenomenon.} The $W^\pm$ and $Z^0$ gauge bosons of the weak interaction have masses $\sim 80$ and $\sim 91$ GeV/$c^2$, with mass ratio $m_W/m_Z \approx 0.882$ and on-shell Weinberg angle $\sin^2\theta_W = 1 - m_W^2/m_Z^2 \approx 0.223$. The weak force has finite range $r_W \sim 10^{-18}$ m.

\textbf{Substrate primitives invoked.}

\begin{itemize}
    \item \textbf{Cosserat couple-stress modulus} $\gamma_c$ (Ch.~\ref{ch:vol9_mechanical_characteristics}~$\S$\ref{sec:vol9_mech_couple_stress}; Ch.~\ref{ch:vol9_magnetic_microrotational_characteristics}~$\S$\ref{sec:vol9_weak_force_range_wz_cutoff}); substrate-native microrotational bending inductance setting the bond-level resistance to differential micro-rotation.
    \item \textbf{Cosserat characteristic length} $l_c = \sqrt{\gamma_c/G_{vac}}$ (Ch.~\ref{ch:vol9_magnetic_microrotational_characteristics}~$\S$\ref{sec:vol9_weak_force_range_wz_cutoff}); identified with the weak-force range $r_W$.
    \item \textbf{Vacuum Poisson ratio} $\nu_{vac} = 2/7$ (Ch.~\ref{ch:vol9_mechanical_characteristics}~$\S$\ref{sec:vol9_mech_poisson}); Ax~1 LC + Ax~2 $\alpha$ interlink, NOT a free parameter.
    \item \textbf{Perpendicular Axis Theorem at bond-cylinder diameter $d = \ell_{node}$} (Ax~1 hard-sphere exclusion; Ch.~\ref{ch:vol9_magnetic_microrotational_characteristics}~$\S$\ref{sec:vol9_weak_force_range_wz_cutoff}): fixes $J = 2I$ for a uniform cross-section.
    \item \textbf{Below-cutoff evanescent-wave mathematics on the microrotational sector} (Ch.~\ref{ch:vol9_magnetic_microrotational_characteristics}~$\S$\ref{sec:vol9_weak_force_range_wz_cutoff}; \kbleaf{ave-kb/vol2/particle-physics/ch05-electroweak-mechanics/gauge-boson-masses.md} line~39): static field equation becomes massive Helmholtz $\nabla^2\theta - \theta/l_c^2 = 0$ with Yukawa Green's function $V \propto e^{-r/l_c}/r$.
\end{itemize}

\textbf{Cross-chapter chain.} Weak interactions on the substrate lack the kinetic energy required to overcome the ambient Cosserat micro-rotation inductance and become evanescent waves at characteristic decay length $l_c = \sqrt{\gamma_c/G_{vac}}$. The $W^\pm$ and $Z^0$ gauge bosons are the substrate's evanescent cutoff modes in two distinct geometric channels (Ch.~\ref{ch:vol9_magnetic_microrotational_characteristics}~$\S$\ref{sec:vol9_weak_force_range_wz_cutoff}; \kbleaf{gauge-boson-masses.md}): $W^\pm$ is the pure longitudinal-torsional evanescent mode with characteristic wavenumber $k \propto G_{vac} J$ (torsional stiffness $\times$ polar moment); $Z^0$ is the transverse-bending evanescent mode with $k \propto E_{vac} I$ (bending stiffness $\times$ second moment of area). The Perpendicular Axis Theorem at a uniform bond cylinder of diameter $d = \ell_{node}$ fixes $J = 2I$, so the mass ratio reduces to $m_W/m_Z = \sqrt{2 G_{vac}/E_{vac}}$. Applying the standard isotropic-elastic identity $E = 2G(1+\nu)$ at the substrate Poisson ratio $\nu_{vac} = 2/7$ (Ch.~\ref{ch:vol9_mechanical_characteristics}~$\S$\ref{sec:vol9_mech_poisson}):

\[
m_W/m_Z = \sqrt{2G/[2G(1+2/7)]} = \sqrt{7/9} = \sqrt{7}/3 \approx 0.882,
\]

and the on-shell Weinberg angle $\sin^2\theta_W = 1 - m_W^2/m_Z^2 = 1 - 7/9 = 2/9 \approx 0.222$ follows directly (canonical-source \kbleaf{src/ave/core/constants.py}). The substrate's electroweak masses are therefore derived from one Cosserat couple-stress modulus $\gamma_c$ + one Cosserat bulk modulus $G_{vac}$ + the substrate-axiomatic Poisson ratio $\nu_{vac} = 2/7$. The EE projection (Ch.~\ref{ch:vol9_magnetic_microrotational_characteristics}~$\S$\ref{sec:vol9_magnetic_ee_translation}, means-test row~10): the same below-cutoff coupling mathematics that governs weak interactions also governs below-cutoff coupling between two LC tanks separated by a transformer leakage inductance --- $\gamma_c$ is the substrate-bond-scale transformer-leakage-inductance characteristic length.

\textbf{Joint kill-switch.} The same $\gamma_c$ Cosserat couple-stress modulus that sets the weak-force range $r_W = l_c$ also sets the B-mode mass-gap $\omega_m^2 = 4 G_c/I_\omega$ that produces $\delta_{strain}$ at $T_{CMB}$ (\S\ref{sec:vol9_app_delta_strain_cosmic_tcc} below). Detection of a stable, freely propagating right-handed neutrino falsifies the $\tfrac{1}{3}G_{vac}$ microrotational boundary condition, destroying the weak-force derivation AND the $\delta_{strain}$ Cosserat-Curie mechanism jointly (Ch.~\ref{ch:vol9_falsification_tests}~$\S$\ref{sec:vol9_neutrino_kill_switch}; \kbleaf{ave-kb/vol4/falsification/ch11-experimental-bench-falsification/epistemology-of-falsification.md}).

\textbf{Discipline class.} Per \texttt{ave-discrimination-check}: the SM Standard Model derives the W/Z masses from the Higgs vacuum-expectation-value $v$ + the electroweak couplings $g$, $g'$, with $m_W = gv/2$, $m_Z = v\sqrt{g^2+g'^2}/2$; the substrate version reaches identical numerical predictions from $\gamma_c$, $G_{vac}$, $\nu_{vac} = 2/7$, and the Perpendicular Axis Theorem. The mass ratio $m_W/m_Z = \sqrt{7}/3 \approx 0.882$ is a substrate-derived consequence of the $\nu_{vac} = 2/7$ Ax~1 LC + Ax~2 $\alpha$ interlink and the Cosserat micropolar structure of the substrate; the SM counterfactual produces the same numerical ratio via Higgs vev mechanics but does not derive the $2/7$ from substrate axioms. The right-handed-neutrino joint-falsifier is substrate-distinct (the SM model can accommodate either the existence or non-existence of right-handed neutrinos as a model parameter without internal contradiction; the substrate model cannot).

\section{Example: Born Rule as Boundary-Joule Extraction at a $\Gamma$ Port}
\label{sec:vol9_app_born_rule_gamma_extraction}

\textbf{Observed phenomenon.} The Born rule $\Pr \propto |\psi|^2$ --- the squared-amplitude scaling of measurement-outcome probability in quantum mechanics, equivalent to a $p = 2$ exponent in the click-rate-versus-substrate-amplitude relation at a detector.

\textbf{Substrate primitives invoked.}

\begin{itemize}
    \item \textbf{Substrate boundary semantics} via $\Gamma = (Z_2 - Z_1)/(Z_2 + Z_1)$ (Op3) at every internal impedance discontinuity (Ch.~\ref{ch:vol9_pin_port_configuration}~$\S$\ref{sec:vol9_substrate_boundary_semantics}).
    \item \textbf{Axiom 3 minimum-reflection} as the substrate's native action principle (Ch.~\ref{ch:vol9_pin_port_configuration}~$\S$\ref{sec:vol9_substrate_boundary_semantics}).
    \item \textbf{Op17 power transmission} $T^2 = 1 - \Gamma^2$ at every substrate port (Ch.~\ref{ch:vol9_pin_port_configuration}~$\S$\ref{sec:vol9_operator_port_behavior}; canonical-source \kbleaf{ave-kb/common/operators.md} Op17 row).
    \item \textbf{Ohmic boundary kinematics} $V^2/Z_{det}$ at the substrate-detector boundary (Ch.~\ref{ch:vol9_dc_electrical_characteristics}~$\S$\ref{subsec:vol9_dc_impedance_trio} bond TLM $Z_0$; Ch.~\ref{ch:vol9_pin_port_configuration}~$\S$\ref{sec:vol9_engineering_translation} EE translation table).
    \item \textbf{Substrate operating-point amplitude statistics} on the Axiom~4 kernel (Ch.~\ref{ch:vol9_saturation_characteristics}~$\S$\ref{sec:vol9_dielectric_specialization}); quadratic-Lagrangian per-site amplitude distribution + cumulant truncation (\kbleaf{ave-kb/vol1/dynamics/ch3-quantum-signal-dynamics/ohmic-decoherence-born.md}).
\end{itemize}

\textbf{Cross-chapter chain.} The substrate's measurement event at a detector is a substrate boundary-Joule extraction event: substrate energy is drawn out of the bulk at a $\Gamma$-port boundary via $V^2/Z_{det}$ Joule kinematics. Substrate amplitude statistics at the boundary are governed by the per-site quadratic-Lagrangian distribution + cumulant truncation; the extraction rate is the boundary-Joule integral over those amplitude statistics; the leading dependence on substrate amplitude is quadratic. The derivation (\kbleaf{ave-kb/vol1/dynamics/ch3-quantum-signal-dynamics/ohmic-decoherence-born.md}) closes the master-equation derivation path end-to-end: master vacuum equation $\to$ boundary-impedance thermalization (FDT) at extraction lattice site $\to$ Joule extraction kinematics $\to$ cumulant truncation under quadratic-Lagrangian per-site amplitude statistics $\to$ quadratic-in-amplitude boundary-Joule extraction-rate scaling.

The $p = 2$ exponent is therefore substrate-derived rather than postulated: $\Pr \propto |\psi|^2$ is the substrate's quadratic-in-amplitude boundary-Joule extraction-rate scaling at every $\Gamma$-port detector boundary. Per the canonical EE translation (\kbleaf{translation-circuit.md} \S 4; \kbleaf{translation-qm.md} Section B): the substrate amplitude $\psi$ is the boundary substrate field; $|\psi|^2$ is the substrate energy density per site; $V^2/Z_{det}$ is the Ohmic extraction rate per Joule kinematics; the $p = 2$ click-rate scaling is the substrate's quadratic-Lagrangian Gaussian-amplitude consequence. The SM-vocabulary measurement-postulate is identified as a substrate-derived consequence, not an independent rule.

\textbf{Discipline class.} In substrate-native terms: the standard-physics ``Born rule $p=2$'' label is parenthetical translation; the substrate-native form is ``quadratic-in-amplitude boundary-Joule extraction-rate scaling at a $\Gamma$-port detector boundary.'' Per \texttt{consistency-vs-emergence} v1.3: this is a Class~D emergence-class derivation --- the $p = 2$ exponent is derived from substrate primitives (master equation + boundary kinematics + amplitude statistics) without postulating a measurement rule. Per \texttt{ave-discrimination-check}: the SM-counterfactual ``Born rule is a postulate'' is consistent with all current observations because no measurement directly tests the substrate's $V^2/Z_{det}$ extraction kinematics at the substrate boundary; the substrate version reaches the same observable predictions plus the substrate-distinct forward predictions (Rice-formula crossing rate, Wald first-passage form, $\hbar\omega$ resonant-mode energy quantum per click).

\section{Example: $\delta_{strain}$ as Cosmic-Scale Temperature Coefficient}
\label{sec:vol9_app_delta_strain_cosmic_tcc}

\textbf{Observed phenomenon.} CODATA $\alpha^{-1} = 137.035999$ sits below the cold-lattice substrate asymptote $\alpha^{-1}_{\text{ideal}} = 4\pi^3 + \pi^2 + \pi \approx 137.0363038$ by the fractional vacuum strain coefficient $\delta_{strain} \approx \num{2.22e-6}$.

\textbf{Substrate primitives invoked.}

\begin{itemize}
    \item \textbf{Substrate bipartite thermal-mode structure} (Ch.~\ref{ch:vol9_temperature_characteristics}~$\S$\ref{subsec:vol9_bipartite_modes}): three translational E-DOFs (gapless acoustic, thermally populated at any $T > 0$) + three microrotational B-DOFs (Cosserat-rotation-sector mass-gap $\omega_m^2 = 4G_c/I_\omega$, thermally frozen below $\hbar\omega_m/k_B$).
    \item \textbf{CMB thermal bath at $T_{CMB} = \qty{2.725}{\kelvin}$} (Ch.~\ref{ch:vol9_temperature_characteristics}~$\S$\ref{subsec:vol9_thermal_population}); Boltzmann factor $\exp(-\hbar\omega_m/k_BT_{CMB}) \sim \exp(-\num{4.3e9}) \approx 0$ for B-modes.
    \item \textbf{Cosserat couple-stress modulus} $\gamma_c$ (Ch.~\ref{ch:vol9_mechanical_characteristics}~$\S$\ref{sec:vol9_mech_couple_stress}; Ch.~\ref{ch:vol9_magnetic_microrotational_characteristics}~$\S$\ref{sec:vol9_rotation_sector_mass_gap}); the shared substrate primitive behind both the B-mode mass-gap AND the weak-force range $l_c$.
    \item \textbf{Asymmetric thermal-mode population ASYM-class scaling} (Ch.~\ref{ch:vol9_temperature_characteristics}~$\S$\ref{subsec:vol9_asym_sym_breaking}; Ch.~\ref{ch:vol9_pin_port_configuration}~$\S$\ref{sec:vol9_port_class_taxonomy} ASYM-N subclasses): $\varepsilon_{eff}(T_{CMB}) = \varepsilon_0(1 - \eta_\varepsilon)$, $\mu_{eff}(T_{CMB}) = \mu_0$.
    \item \textbf{Two effective wave speeds} (Ch.~\ref{ch:vol9_ac_electrical_characteristics}~$\S$\ref{sec:vol9_two_wave_speeds}; INVARIANT-S2 Pitfall~\#5): the thermal modulation acts on the $c_{EM}$ branch ($c_{EM} = 1/\sqrt{\mu_{eff}\varepsilon_{eff}}$, which enters $\alpha$), not the $c_{shear}$ branch.
\end{itemize}

\textbf{Cross-chapter chain.} At any operating temperature $T > 0$, the substrate's gapless translational E-modes are thermally populated by Bose-Einstein occupation while the gapped microrotational B-modes are mass-gap-suppressed by $\exp(-\hbar\omega_m/k_BT)$ (Ch.~\ref{ch:vol9_temperature_characteristics}~$\S$\ref{subsec:vol9_bipartite_modes}; Ch.~\ref{ch:vol9_magnetic_microrotational_characteristics}~$\S$\ref{sec:vol9_rotation_sector_mass_gap}). At $T = T_{CMB}$, the B-mode Boltzmann suppression is total ($\sim e^{-\num{4e9}}$), so $\mu_{eff} = \mu_0$ at cold-lattice while $\varepsilon_{eff} = \varepsilon_0(1 - \eta_\varepsilon)$ is modulated by thermal-photon-bath occupation of E-modes. The asymmetric scaling voids the SYM-class $\alpha$-invariance condition (Ch.~\ref{ch:vol9_saturation_characteristics}~$\S$\ref{sec:vol9_alpha_invariance}; Ch.~\ref{ch:vol9_pin_port_configuration}~$\S$\ref{sec:vol9_port_class_taxonomy}): substituting into $\alpha = e^2/(4\pi\varepsilon_{eff}\hbar c_{EM,eff})$ with $c_{EM,eff} = 1/\sqrt{\mu_0\varepsilon_{eff}}$, the leading-order correction is

\[
\frac{\alpha_{eff}}{\alpha_0} \approx \frac{1}{(1-\eta_\varepsilon)^{1/2}} \approx 1 + \frac{\eta_\varepsilon}{2}
\qquad \Rightarrow \qquad
\delta_{strain} \;\equiv\; -\frac{\delta\alpha^{-1}}{\alpha^{-1}_0} \approx \frac{\eta_\varepsilon}{2}.
\]

With CODATA $\alpha^{-1} = 137.035999$ and the cold-lattice substrate asymptote $\alpha^{-1}_{\text{ideal}} = 4\pi^3 + \pi^2 + \pi \approx 137.0363038$:

\[
\delta_{strain}(T_{CMB}) \approx \num{2.22e-6}, \qquad \eta_\varepsilon(T_{CMB}) \approx \num{4.45e-6}.
\]

The substrate's cosmic-scale temperature coefficient $\delta_{strain}$ is therefore the same Cosserat-Curie mechanism that produces engineering-scale TCC entries in component datasheets, applied at cosmic scale with the CMB photon bath as the thermal driver (\kbleaf{ave-kb/vol3/cosmology/ch05-dark-sector/delta-strain-cosmic-tcc.md}; mirrored at the vol3 translation-circuit \S 9 EE TCR/TCC engineering-correction row).

\textbf{Forward prediction (substrate-distinct).} The substrate predicts roughly linear cosmic-temperature $\alpha$-drift for $T_{CMB} \ll T \ll T_{B-gap} \sim 10^{10}$~K (\kbleaf{ave-kb/vol3/cosmology/ch05-dark-sector/delta-strain-cosmic-tcc.md} \S6). Quasar absorption-line $\Delta\alpha/\alpha$ measurements at higher redshift access this thermal regime; current bounds $\sim 10^{-5}$ are consistent with the substrate-predicted scaling at the back-subtracted $\eta_\varepsilon$ (Ch.~\ref{ch:vol9_falsification_tests}~$\S$\ref{subsec:vol9_quasar_alpha_null}). The substrate-distinct prediction beyond standard SM/QED null is the positive thermal scaling (SIGN) at higher cosmic $T$; the magnitude $\eta_\varepsilon(T)$ is not derived from substrate statistical mechanics (the candidate thermal-occupation derivation undershoots by $\sim$31 orders of magnitude and is not substrate-distinct), so the magnitude is a definitional residual --- the AVE-distinct handle is the direction, not the magnitude (Ch.~\ref{ch:vol9_temperature_characteristics}~$\S$\ref{sec:vol9_temperature_open}).

\textbf{Discipline class.} Per \texttt{consistency-vs-emergence} v1.3: substrate-mechanism class is Class~B, predicting the SIGN (\kbleaf{ave-kb/vol3/cosmology/ch05-dark-sector/delta-strain-cosmic-tcc.md}); the magnitude $\eta_\varepsilon(T_{CMB})$ is not derived from substrate statistical mechanics (the candidate thermal-occupation derivation undershoots by $\sim$31 orders of magnitude and is not substrate-distinct). Per \texttt{ave-discrimination-check}: the SM-counterfactual --- standard QED has no thermal $\alpha$-running mechanism at cosmic photon-bath occupation; CODATA $\alpha^{-1}$ is treated as a fundamental empirical constant with no substrate-mechanism explanation for the precise value. The substrate version identifies $\alpha^{-1}_{\text{CODATA}} = \alpha^{-1}_{\text{ideal}}(1 - \delta_{strain})$ with the asymptote ($4\pi^3 + \pi^2 + \pi$) a named geometric identification and the correction ($\delta_{strain}$ at $T_{CMB}$) a \textbf{definitional residual} ($1-$CODATA$/\alpha_{cold}$, NOT derived from substrate primitives --- only its SIGN is substrate-set). The joint kill-switch (right-handed neutrino observation) ties this prediction to the W/Z derivation of \S\ref{sec:vol9_app_wz_transformer_leakage} via the shared $\gamma_c$ primitive.

\section{Example: The Electron Node as a Two-Domain N-Port}
\label{sec:vol9_app_node_2domain_nport}

\textbf{Observed phenomenon.} The electron presents one external electromagnetic coupling (it radiates and absorbs photons at the matched vacuum impedance) while its mass and charge are \emph{confined} --- not radiated. A faithful equivalent circuit must reproduce this asymmetry: \emph{one} external port plus \emph{two} confined internal channels.

\textbf{Substrate primitives invoked.}

\begin{itemize}
    \item \textbf{Three wired reactance channels} (Ch.~\ref{ch:vol9_pin_port_configuration} \S Device Circuit Models; canonical at \kbleaf{device-circuit-models.md}:141--149): the EM channel $Z_{EM} \equiv Z_0 = \sqrt{\mu_0/\varepsilon_0} \approx 376.73\,\Omega$ ($\Gamma_{EM} = 0$, matched), the shear channel $Z_{shear} = \rho_{bulk}\, c_{shear}$ ($\Gamma_{shear} \to -1$, confined --- the CHARGE-``3''), and the bulk channel $Z_{bulk} = \sqrt{2}\,\rho_{bulk}\, c_0$ at $K = 2G$ ($\Gamma_{bulk} \to -1$, confined --- the MASS-``3'').
    \item \textbf{Mixed impedance domains} (Ch.~\ref{ch:vol9_mechanical_characteristics} \S Node-model channel tags): $Z_{EM}$ is an \emph{electrical} impedance ($\Omega$); $Z_{shear}$ and $Z_{bulk}$ are \emph{mechanical/acoustic} impedances ($\rho \times$ speed, Pa$\cdot$s/m), $\sim$12 OOM off $Z_0$ and a unit change away.
    \item \textbf{The EM-$\Omega$ $\leftrightarrow$ mechanical-$\rho c$ TKI-transformer} (Ch.~\ref{ch:vol9_topological_characteristics} $\xi_{topo}$; \kbleaf{device-circuit-models.md}:205--210; the transducer node is provisional, not yet ratified): two ports in different units cannot be directly wired --- a domain crossing requires a TRANSDUCER (an electro-mechanical change-of-reference), realized as an ideal transformer (turns-ratio scaling), not a direct wire and not a gyrator.
    \item \textbf{$\alpha$ localized to the transducer} via the critical packing fraction $p_c = 8\pi\alpha$ (\kbleaf{src/ave/core/constants.py}, \texttt{P\_C}): the $\alpha$-echo localizes to \emph{exactly} the EM$\leftrightarrow$mechanical transformer turns-ratio (Ch.~\ref{ch:vol9_breakdown_characteristics} dielectric-rupture consistency identity).
    \item \textbf{Fork-A conservative shear$\leftrightarrow$bulk coupling} (Ch.~\ref{ch:vol9_magnetic_microrotational_characteristics} \S\ref{sec:vol9_magnetic_forka}): a bounded, lossless, helicity-transferring skew-Hermitian coupling between the two confined mechanical channels (PARTIAL --- the form is forced, the non-reciprocity magnitude is imposed).
\end{itemize}

\textbf{Cross-chapter chain.} Draw the electron as a \textbf{2-domain N-port}: the single EM port ($Z_0 = 376.73\,\Omega$, $\Gamma_{EM} = 0$, the sole external channel --- how the interior couples to the far field) wired through an explicit ideal transformer into the mechanical domain, where the two confined channels live --- the bulk channel carrying the MASS-``3'' ($Z_{bulk} = \sqrt{2}\,\rho_{bulk}\, c_0$ at $K = 2G$) and the shear channel carrying the CHARGE-``3'' ($Z_{shear} = \rho_{bulk}\, c_{shear}$), both terminating at the confinement surface where $\Gamma \to -1$ (the driver \kbleaf{src/scripts/vol\_9\_device/node\_2domain\_nport.py} implements this as an explicit ideal-transformer two-port). The transducer carries $\alpha$: the lumped-($\Omega$, per-cell) $\to$ specific-($\rho c$) reconciliation runs through $Z_{referred} = Z_{EM}\,\xi_{topo}^2/(p_c\,\ell_{node}^2)$ with $p_c = 8\pi\alpha$, and the $(\xi_{topo}^2\mu_0)/(p_c\,\ell_{node}^2) = \rho_{bulk}$ identity makes the $\alpha$-echo \emph{localize visibly} to that single transformer factor (\kbleaf{device-circuit-models.md}:205--210). The two confined channels couple to each other --- but \emph{only} through a conserved Hamiltonian pair (the Fork-A skew-circulator; Ch.~\ref{ch:vol9_magnetic_microrotational_characteristics} \S\ref{sec:vol9_magnetic_forka}), never through a shared phasor ($A1 \perp T2$).

\textbf{Discipline class.} Per \texttt{consistency-vs-emergence} v1.3: the 2-domain N-port is a \textbf{Class~C CONSISTENCY re-expression} of the already-canonical three-impedance law as a wired equivalent-circuit MODEL --- it originates no new substrate primitive (\kbleaf{device-circuit-models.md}:127--131). The transducer node is provisional: it transduces units losslessly but does not supply a derived coupling \emph{mechanism}. $\alpha = p_c/8\pi$ at the transducer is a \emph{consistency identity}, NOT an $\alpha$ derivation; and the shear$\leftrightarrow$bulk coupling is partial: the form is forced but the non-reciprocity magnitude is imported, not derived (\kbleaf{common/form-deriving-value-importing.md}). Per \texttt{ave-discrimination-check}: the SM-counterfactual draws the electron as a point charge with a fixed coupling $\alpha$ --- the substrate version reaches the same external behavior (one matched radiative port) but exhibits \emph{where} the coupling lives (the transducer turns-ratio) rather than positing it.

\section{Application Information: Worked Cell Analysis (SPICE Phase-1 Ladder)}
\label{sec:vol9_app_worked_spice_ladder}

This section is the datasheet's \textbf{worked circuit analysis} --- the vacuum cell's constitutive models exercised through a real SPICE engine (\texttt{ngspice-46}), each rung stated as an analytic target, a measured value, and a recovery error. It is a consistency-class validation ladder (a validate-on-known, per \texttt{consistency-vs-emergence} v1.3): the linear rungs are known-analytic calibrations; the Axiom-4 rungs are manifestation cross-checks (the \texttt{.lib} behavioral source \emph{is} the canonical kernel, not a fit). The netlists are committed at \kbleaf{src/scripts/vol\_4\_engineering/spice\_ladder\_artifacts/}; the full result record is \kbleaf{research/2026-07-04\_spice-phase1-ladder\_result.md}. The cells are drawn in Fig.~\ref{fig:vol9_spice_ladder_cells}.

% spice_ladder_cells.tex — SPICE phase-1 validation-ladder cells (Vol 9 Ch.13)
% Renders the committed netlists in src/scripts/vol_4_engineering/spice_ladder_artifacts/.
% House style: WHITE, american circuitikz, reference designators mirroring the .cir.
\begin{figure}[htbp]
\centering
\begin{circuitikz}[american, scale=0.82, transform shape]
    % ---- (a) RC charging cell (rung 1a) ----
    \node[font=\bfseries\scriptsize] at (2.0, 2.3) {(a) RC transient — rung 1a};
    \draw (0,0) node[left, font=\scriptsize] {$V_1$}
        to[V, invert] (0,-2.2);       % step source
    \draw (0,0) to[R, l=$R_1$, a=$1\,\mathrm{k}\Omega$] (2.5,0);
    \draw (2.5,0) to[C, l=$C_1$, a=$1\,\mu\mathrm{F}$, *-] (2.5,-2.2)
        node[ground] {};
    \draw (0,-2.2) -- (2.5,-2.2);
    \node[circle, fill=black, inner sep=1pt] at (2.5,0) {};
    \node[right, font=\scriptsize] at (2.7,0) {$v(N_C)$};
    \node[font=\scriptsize, align=center] at (1.25,-3.0)
        {$\tau = R_1 C_1$};

    % ---- (b) LC oscillator cell (rung 1b) ----
    \begin{scope}[xshift=6.6cm]
    \node[font=\bfseries\scriptsize] at (1.6, 2.3) {(b) LC tank — rung 1b};
    \draw (0,0) to[C, l=$C_1$, a=$1\,\mu\mathrm{F}$] (3.2,0);
    \draw (0,0) to[L, l_=$L_1$, a^=$1\,\mathrm{mH}$] (0,-2.2);
    \draw (3.2,0) -- (3.2,-2.2);
    \draw (0,-2.2) -- (3.2,-2.2) node[ground, midway] {};
    \node[circle, fill=black, inner sep=1pt] at (3.2,0) {};
    \node[right, font=\scriptsize] at (3.35,0) {$v(N_C)$};
    \node[font=\scriptsize, align=center] at (1.6,-3.0)
        {$f_{\mathrm{res}} = \dfrac{1}{2\pi\sqrt{L_1 C_1}}$, \; $\mathrm{IC}{=}1\,\mathrm{V}$};
    \end{scope}

    % ---- (c) 1D LC ladder (rungs 4 + 5) ----
    \begin{scope}[yshift=-5.4cm]
    \node[font=\bfseries\scriptsize] at (3.4, 1.5) {(c) $N$-cell LC ladder — rungs 4 (dispersion) \& 5 (biased small-signal)};
    \draw (0,0) node[left, font=\scriptsize] {AC in}
        to[L, l=$L$] (2,0)
        to[L, l=$L$] (4,0)
        to[L, l=$L$] (6,0)
        to[short] (7,0)
        node[right, font=\scriptsize] {$\cdots\;Z_c$};
    \draw (2,0) to[C, l=$C_{\mathrm{eff}}$, *-] (2,-1.8) node[ground] {};
    \draw (4,0) to[C, l=$C_{\mathrm{eff}}$, *-] (4,-1.8) node[ground] {};
    \draw (6,0) to[C, l=$C_{\mathrm{eff}}$, *-] (6,-1.8) node[ground] {};
    \node[font=\scriptsize, align=center] at (3.4,-2.6)
        {rung 4: $C_{\mathrm{eff}}{=}C_0$ (cold); rung 5: $C_{\mathrm{eff}}{=}C_0/S(V_b)^3$ (biased small-signal, A1 varactor)};
    \node[font=\scriptsize, align=center] at (3.4,-3.3)
        {band: $\;ka = 2\arcsin(\omega/2\omega_0)$, \; $\omega_0 = 1/\sqrt{LC_{\mathrm{eff}}}$};
    \end{scope}
\end{circuitikz}
\caption[SPICE phase-1 validation-ladder cells]{The SPICE phase-1 validation-ladder cells, drawn from the committed netlists (\kbleaf{src/scripts/vol\_4\_engineering/spice\_ladder\_artifacts/}). (a) The RC charging cell (\texttt{spice\_ladder\_rung1\_rc.cir}); (b) the pre-charged LC tank (\texttt{spice\_ladder\_rung1\_lc.cir}); (c) the $N$-cell LC ladder used for the dispersion band (rung 4, cold $C_{\mathrm{eff}}{=}C_0$) and the biased small-signal shift (rung 5, $C_{\mathrm{eff}}{=}C_0/S(V_b)^3$). Reference designators mirror the \texttt{.cir} files. All quantities are lumped-network (matching coordinates throughout, A46-clean).}
\label{fig:vol9_spice_ladder_cells}
\end{figure}


\subsection{RC / LC transients ($\mathtt{.tran}$)}
\label{subsec:vol9_app_spice_rc_lc}

\textbf{(a) RC charging.} A $1\,\mathrm{V}$ step drives $R_1 = 1\,\mathrm{k}\Omega$ in series with $C_1 = 1\,\mu\mathrm{F}$ to ground (Fig.~\ref{fig:vol9_spice_ladder_cells}a; \texttt{spice\_ladder\_rung1\_rc.cir}). \emph{Symbolic:} $v_C(t) = V_0\,(1 - e^{-t/\tau})$, $\tau = R_1 C_1$; the $63.2\%$ ($1 - 1/e$) crossing recovers $\tau$. \emph{Numeric:} $\tau_{\mathrm{analytic}} = R_1 C_1 = \qty{1.000000e-3}{\second}$; \texttt{ngspice} measured $\qty{1.000000e-3}{\second}$ (rel-error $\num{1.21e-7}$, tol $\num{2e-2}$).

\textbf{(b) LC oscillation.} An undamped tank $L_1 = 1\,\mathrm{mH}$, $C_1 = 1\,\mu\mathrm{F}$, cap pre-charged to $1\,\mathrm{V}$ (\texttt{IC}), rings (Fig.~\ref{fig:vol9_spice_ladder_cells}b; \texttt{spice\_ladder\_rung1\_lc.cir}). \emph{Symbolic:} $f_{\mathrm{res}} = 1/(2\pi\sqrt{L_1 C_1})$; recovered from the Hann-windowed FFT peak. \emph{Numeric:} $f_{\mathrm{res,analytic}} = \qty{5032.92}{\hertz}$; \texttt{ngspice} measured $\qty{5039.67}{\hertz}$ (rel-error $\num{1.34e-3}$, tol $\num{2e-2}$). This is the lumped-network analog of the node LC tank whose element values are $L_{\mathrm{cell}} = \mu_0\ell_{node}$, $C_{\mathrm{cell}} = \varepsilon_0\ell_{node}$ (Ch.~\ref{ch:vol9_dc_electrical_characteristics}; \kbleaf{src/ave/core/constants.py} \texttt{L\_CELL}, \texttt{C\_CELL}).

\subsection{Statics solve ($\mathtt{.OP}$)}
\label{subsec:vol9_app_spice_op}

A 24-node / 32-edge random resistor graph with $1\,\mathrm{mA}$ injected at the far node and node~0 grounded (\texttt{spice\_ladder\_rung3\_poisson.cir}) is solved three ways. \emph{Symbolic:} the SPICE \texttt{.OP} MNA system is the weighted graph-Laplacian $L\,\mathbf{v} = \mathbf{i}$ with the ground row/column deleted (the neutrality boundary condition). \emph{Numeric:} $\max|v_{\mathrm{ngspice}} - v_{\mathrm{MNA}}| = \qty{4.17e-10}{\volt}$ and $\max|v_{\mathrm{ngspice}} - v_{\mathrm{Laplacian}}| = \qty{4.17e-10}{\volt}$ (the residual tracks the \texttt{print} precision --- a text-serialization artifact, not a solver difference; tol $\qty{1e-8}{\volt}$). The real-engine \texttt{.OP} MNA system \emph{is} the srs statics Laplacian.

\subsection{Lattice dispersion ($\mathtt{.AC}$, cold ladder)}
\label{subsec:vol9_app_spice_ac_dispersion}

A 40-cell LC ladder ($L = 1\,\mu\mathrm{H}$, $C = 1\,\mathrm{nF}$, so $\omega_0 = 1/\sqrt{LC} = \qty{3.16e7}{\radian\per\second}$), matched-terminated in $Z_c = \sqrt{L/C}$, is driven by a $1\,\mathrm{V}$ \texttt{.AC} source (Fig.~\ref{fig:vol9_spice_ladder_cells}c, cold; \texttt{spice\_ladder\_rung4\_dispersion.cir}). \emph{Symbolic:} the discrete-lattice band is $\omega(k) = 2\omega_0\,|\sin(ka/2)|$, i.e.\ $ka = 2\arcsin(\omega/2\omega_0)$, a band with a zone-edge cutoff at $\omega_{\max} = 2\omega_0$ ($ka = \pi$) --- \emph{not} the continuum $\omega = ck$ line. The phase advance per cell $\Delta\varphi = \arg(V_{m+1}/V_m)$ is fit over an interior window. \emph{Numeric} (pass-band $0.05 < ka < 0.9\pi$): median rel-error in $ka$ $\num{1.87e-3}$, max $\num{3.03e-2}$ (tol $\num{5e-2}$). The measured band follows the $\sin$-bending away from the linear continuum line --- the substrate-native lattice transmission line, not a continuum wire.

\subsection{Bias couples to the wave ($\mathtt{.AC}$, biased small-signal)}
\label{subsec:vol9_app_spice_ac_biased}

The same 40-cell ladder, now with the Axiom-4 metric varactor ($C_{\mathrm{eff}} = C_0/S(V)$, A1-divergent, keyed on $V_{\mathrm{snap}}$) as the shunt element, is DC-biased and re-swept (Fig.~\ref{fig:vol9_spice_ladder_cells}c, biased; \texttt{spice\_ladder\_rung5\_bias\_hi.cir} / \texttt{\_lo.cir}). This is the datasheet's \textbf{large-signal-vs-small-signal} distinction made concrete: the large-signal chord varactor is $C_{\mathrm{eff}} = C_0/S(V)$, but the \texttt{.AC} propagation is governed by the \emph{small-signal differential} capacitance
\begin{equation}
C_{\mathrm{ss}} = \left.\frac{\mathrm{d}Q}{\mathrm{d}V}\right|_{V_b} = \frac{C_0}{S(V_b)^3},
\qquad Q = \frac{C_0 V}{S(V)},\ \ S(V) = \sqrt{1 - (V/V_{\mathrm{snap}})^2},
\label{eq:vol9_smallsignal_ceff}
\end{equation}
(verified in \texttt{ngspice} against the analytic $\mathrm{d}Q/\mathrm{d}V$ to $\sim\!\num{1e-10}$). \emph{Symbolic:} biasing raises $C_{\mathrm{eff}}$, drops $\omega_0(V_b) = 1/\sqrt{L\,C_{\mathrm{ss}}}$, and at fixed $\omega$ the phase-per-cell $ka$ \emph{increases} (the wave slows); the S(A)-predicted band is $\omega(k) = 2\omega_0(V_b)\,|\sin(ka/2)|$. \emph{Numeric:} at $V_{\mathrm{hi}} = 0.8\,V_{\mathrm{snap}} = \qty{408799}{\volt}$, $S = 0.6$, so $C_{\mathrm{eff}} = C_0/S^3 = 4.63\,C_0$ and the band drops $2.15\times$; the measured bias-shift $\Delta ka = ka(V_{\mathrm{hi}}) - ka(V_{\mathrm{lo}})$ tracks the S(A) prediction to median rel-error $\num{4.03e-3}$, max $\num{1.76e-2}$ (tol $\num{5e-2}$). This is the smallest in-silico rehearsal of the bias-couples-to-wave (DC$\to$AC) class the pre-registered vacuum-birefringence falsifier lives in --- a consistency rehearsal, \emph{not} the falsifier, not a chord.

\subsection{Measured-vs-analytic summary}
\label{subsec:vol9_app_spice_summary}

\begin{table}[htbp]
\centering\footnotesize
\begin{tabular}{@{}l l l l l@{}}
\toprule
\textbf{Rung / analysis} & \textbf{Analytic target} & \textbf{Measured} & \textbf{Rel.\ error} & \textbf{Class} \\
\midrule
RC $\mathtt{.tran}$ ($\tau$) & $\qty{1.000000e-3}{\second}$ & $\qty{1.000000e-3}{\second}$ & $\num{1.21e-7}$ & consistency \\
LC $\mathtt{.tran}$ ($f_{\mathrm{res}}$) & $\qty{5032.92}{\hertz}$ & $\qty{5039.67}{\hertz}$ & $\num{1.34e-3}$ & consistency \\
A1 varactor $C_0/S$ (kernel) & $1/S(V)$ canonical & $1/S(V)$ ngspice & $\num{3.94e-7}$ & manifestation \\
T2 collapse $C_0 S$ (kernel) & $S(V)$ canonical & $S(V)$ ngspice & $\num{4.83e-8}$ & manifestation \\
Poisson $\mathtt{.OP}$ (24 nodes) & $v_{\mathrm{MNA}}$ & $v_{\mathrm{ngspice}}$ & $\qty{4.17e-10}{\volt}$ & consistency \\
Dispersion $\mathtt{.AC}$ ($ka$) & $2\arcsin(\omega/2\omega_0)$ & $ka$ (phase/cell) & $\num{1.87e-3}$ med & manifestation \\
Biased small-signal ($\Delta ka$) & $\Delta ka$ from $C_0/S^3$ & $\Delta ka$ measured & $\num{4.03e-3}$ med & consistency (DC$\to$AC) \\
\bottomrule
\end{tabular}
\caption{Worked SPICE phase-1 ladder: analytic target vs \texttt{ngspice-46} measurement, per rung. All seven rungs pass their tolerance; the HALT-gate never tripped. Source: \kbleaf{research/2026-07-04\_spice-phase1-ladder\_result.md}; netlists at \kbleaf{src/scripts/vol\_4\_engineering/spice\_ladder\_artifacts/}.}
\label{tab:vol9_spice_ladder_summary}
\end{table}

\textbf{Discipline class.} Per \texttt{consistency-vs-emergence} v1.3: the RC/LC/OP/dispersion rungs are \textbf{consistency-class} (engine integrates a known-analytic circuit) or \textbf{manifestation-class} (the \texttt{.lib} kernel \emph{is} the Axiom-4 kernel). The biased small-signal rung is a \textbf{DC$\to$AC consistency rehearsal} of the falsifier class, not the falsifier and not a chord (\kbleaf{research/2026-07-04\_spice-phase1-ladder\_result.md}). Empirical-driver discipline (A-Rule~10): three \texttt{.lib} ngspice-syntax bugs and two measurement artifacts (a \texttt{.dc}-sweep lag and a bias/2-divider topology bug) fired at first live parse and were named + fixed, not papered over --- the Rule-10 payoff.

\section{Cross-Chapter Citation Map}
\label{sec:vol9_app_cross_chapter_citation_map}

The table below lists each application example and the Vol~9 chapters from which it draws substrate primitives. Each cell entry names the primary substrate primitive(s) the application invokes from that chapter; readers navigating an example may use this table to locate the substrate-physics spec entry directly.

\begin{table}[htbp]
\centering
\renewcommand{\arraystretch}{1.40}
\footnotesize
\setlength{\tabcolsep}{3pt}
\begin{tabular}{@{}p{0.21\linewidth} p{0.075\linewidth} p{0.075\linewidth} p{0.075\linewidth} p{0.075\linewidth} p{0.075\linewidth} p{0.075\linewidth} p{0.075\linewidth} p{0.075\linewidth} p{0.075\linewidth}@{}}
\toprule
\textbf{Application example} & \textbf{Ch.\,3} & \textbf{Ch.\,4} & \textbf{Ch.\,5} & \textbf{Ch.\,6} & \textbf{Ch.\,7} & \textbf{Ch.\,9} & \textbf{Ch.\,10} & \textbf{Ch.\,11} & \textbf{Ch.\,12} \\
\midrule
$\alpha$ derivation (\S\ref{sec:vol9_app_electron_lc_alpha}) & $\Gamma=-1$ TIR (Ax~3) & $L_{cell}$, $C_{cell}$, $Z_0$, $\xi_{topo}$ & --- & --- & $u_0^*$ kernel closure & --- & --- & $(2,3)$ trefoil, $0_1$ unknot, $T \to 2T$ & $u_0^*$ magic angle \\
Schwarzschild + Machian $G$ (\S\ref{sec:vol9_app_gravity_gradient_index}) & SYM-class $\Gamma=0$ ports & $Z_0$, $T_{EM}$ & $c_{EM}$ vs $c_{shear}$ at $A_0$ & --- & SYM-class $\alpha$-invariance & $G_{vac}$, $K=2G$, $\nu=2/7$, $c_{shear}$ & --- & --- & Machian $G = c^4/(7\xi T_{EM})$ (mixed), $u_0^*$ \\
$W$/$Z$ as transformer-leakage cutoff (\S\ref{sec:vol9_app_wz_transformer_leakage}) & --- & --- & --- & --- & --- & $\gamma_c$, $\nu_{vac}=2/7$, PAT $J=2I$ & $l_c$, $m_W/m_Z = \sqrt{7}/3$, $\sin^2\theta_W = 2/9$ & --- & --- \\
Born-rule $p=2$ as boundary-Joule extraction (\S\ref{sec:vol9_app_born_rule_gamma_extraction}) & Op3 $\Gamma$, Ax~3, Op17 $T^2=1-\Gamma^2$, $V^2/Z_{det}$ & $Z_0$ at boundary & --- & --- & Ax~4 amplitude statistics & --- & --- & --- & --- \\
$\delta_{strain}$ as cosmic-scale TCC (\S\ref{sec:vol9_app_delta_strain_cosmic_tcc}) & ASYM-class port subclass & $\varepsilon_0$, $\mu_0$ baseline & $c_{EM}$ branch at thermal load & Cosserat-Curie + $\eta_\varepsilon$ & ASYM $\alpha$-invariance void & --- & $\gamma_c$, $m_\omega$ B-mode gap & --- & --- \\
\bottomrule
\end{tabular}
\renewcommand{\arraystretch}{1.0}
\caption{Cross-chapter citation map: per application example, the substrate primitive(s) drawn from each Vol~9 chapter.}
\label{tab:vol9_app_citation_map}
\end{table}

\noindent
\textbf{Reading the table.} Each row is an application example treated in a section above; each column is a Vol~9 chapter whose spec table or canonical-source routing contributes a substrate primitive to the example. Entries name the specific primitive(s) invoked, not full citation chains --- the canonical leaves are cited in the individual section bodies. Chapters not listed in the header (Ch.~2 absolute maximum ratings; Ch.~8 breakdown characteristics; Ch.~14 phase diagrams; Ch.~15 falsification tests) appear in individual examples as cross-references rather than primitive-source columns:

\begin{itemize}
    \item \textbf{Ch.~\ref{ch:vol9_absolute_maximum_ratings}} (Absolute Maximum Ratings) supplies $V_{snap}$, $V_{yield}$, $E_S$, $B_{snap}$, $T_{melt}$ as rupture-boundary references for the $\alpha$ derivation (electron LC tank at $V_{yield}$ closure) and the Cosserat-Curie mechanism (B-mode gap thermal cutoff at $T_{B-gap}$).
    \item \textbf{Ch.~\ref{ch:vol9_breakdown_characteristics}} (Breakdown) supplies the Regime~IV substrate-rupture framing referenced by the gravity / BH discussion (interior = Regime~IV ruptured plasma per Ch.~\ref{ch:vol9_cosmological_characteristics}~$\S$\ref{sec:vol9_cosmo_substrate_rupture}) and the Schwinger $E_S$ field referenced by the cosmic vs.\ terrestrial null framing (Ch.~\ref{ch:vol9_falsification_tests}~$\S$\ref{subsec:vol9_schwinger_terrestrial_null}).
    \item \textbf{Ch.~\ref{ch:vol9_phase_diagrams}} (Phase Diagrams) supplies the four-regime partition of the operating-point axis ($r = A_0/A_{yield}$) that locates each example on the substrate's overall phase diagram; the $\alpha$ derivation operates at the Regime~II $\to$ III boundary at $r_2 = \sqrt{3}/2$ (spin-2 sector), gravity at deep Regime~I ($r \ll 1$), the BH interior at Regime~IV, and the cosmic-scale $\delta_{strain}$ at low-amplitude thermal-bath ASYM (a sub-regime within Regime~I documented by the ASYM-N subclass enumeration).
    \item \textbf{Ch.~\ref{ch:vol9_falsification_tests}} (Falsification) supplies the kill-switch tests that empirically interrogate the substrate primitives invoked by each example: PONDER-05 for the Ax~4 dielectric specialization (referenced by $\alpha$ derivation through the operating-point modulation and by $\delta_{strain}$ as the bench-scale analog of the thermal-mode ASYM mechanism); HOPF-01 for the $(p,q)$ chiral-resonance shift (referenced by $\alpha$ derivation through the $(2,3)$ topology); right-handed neutrino joint kill-switch for W/Z + $\delta_{strain}$ shared $\gamma_c$; three-route $u_0^*$ convergence for $\alpha$ + Machian $G$ + cosmological $\mathcal{J}_{cosmic}$.
\end{itemize}

\section{Summary}
\label{sec:vol9_app_summary}

Each application example treated above is the substrate's natural response under specified operating conditions, with the substrate primitives in Chs.~2--12 combining to produce the observed engineering quantity. No example introduces a new substrate primitive; no example postulates a physical phenomenon beyond what the substrate axioms (Ax~1 K4 Cosserat lattice + Ax~2 topo-kinematic isomorphism $\xi_{topo}$ + Ax~3 minimum reflection + Ax~4 universal saturation kernel) plus the structural-closure framing (one scale primitive $\ell_{node}$ + one cosmological initial-data parameter $\hat{\Omega}_{freeze}$) already supplies. The framework's predictive content is the joint constraint across application examples: $\alpha$, $G$, $m_W/m_Z$, $\delta_{strain}$, the Born-rule $p=2$ exponent, and (more broadly) every observable in this datasheet derive from one operating point $u_0^*$ + the substrate axioms; if any single derivation chain breaks, the framework breaks.

The substrate is natural; the application examples document the substrate-mechanism behind specific observed phenomena; engineering measures; AVE substrate-physics derives. The next-step datasheet entries continue the same pattern: Ch.~\ref{ch:vol9_phase_diagrams} consolidates the operating-point + cosmic-phase partition; Ch.~\ref{ch:vol9_falsification_tests} catalogs the empirical kill-switches; Ch.~\ref{ch:vol9_cross_volume_reference} is the auto-generated parameter $\to$ derivation map.
